% Simplified requirement command using tcolorbox
% Usage: \requirement{R-ID}{title}{description}{type}{tags}
\newcommand{\requirement}[5]{%
    % Create hyperlink target for the requirement
    \hypertarget{#1}{}%
    
    % Display requirement using tcolorbox
    \noindent
    \par
    \begin{tcolorbox}[
        enhanced,
        breakable,
        rounded corners,
        frame empty,
        colback=gray!7,
        colbacktitle=white,
        coltitle=black,
        fonttitle=\bfseries,
        before skip=3pt,
        after skip=3pt,
        title={
          \hyperlink{source_#1}{\textbf{#1: #2}}
          \hfill
          \scriptsize{\textmd{#4 | #5}}
        },
        overlay={
          \draw[gray, line width=1pt]
              (title.south west) -- (title.south east);
        }
      ]
      #3\par
    \end{tcolorbox}%
}

% Inline requirement reference command
% Usage: \inlineReq{R-ID}
\newcommand{\inlineReq}[1]{%
    \hyperlink{#1}{[\rightarrow #1]}%
    \hypertarget{source_#1}{}
}

% Inline requirement reference command
% Usage: \inlineReq{R-ID}
\newcommand{\reqRef}[1]{%
    \hyperlink{#1}{[#1]}%
}
% ── TODO / Comment commands ──────────────────────────────────────────
% Annotations now handled by fixme package (see packages.tex).
% Use: \fxnote{}, \fxwarning{}, \fxerror{}, \fxfatal{}
% These work inside floats, captions, and any other context.

% Fix for Biblatex IEEE style missing macro
\renewbibmacro*{volume+number+eid}{%
  \printfield{volume}%
  \setunit*{\addnbthinspace}%
  \printfield{number}%
  \setunit{\addcomma\space}%
  \printfield{eid}}
